% GENERATED FILE. Edit canonical lessons, then run npm run generate:books.
% canonical-source-hash: fnv1a64:febc0aceaf1793bf
% canonical-lessons: UR-C515-tasvir, UR-C515-musavvir, UR-C515-khaka, UR-C515-shahkar, UR-C515-jhalakna

\chapter{Pictures and Painters}
\label{ch:ur-pictures-and-painters}
\hlchaptermodality{515}

\begin{chapteropening}
\textbf{By the end of this chapter:} \emph{I can talk about pictures and the people who paint them.}

Complete the last lesson of chapter 515: I can talk about pictures and the people who paint them.
\end{chapteropening}

\section[\texorpdfstring{\ur{تصویر} (tasvīr)}{تصویر (tasvīr)}]{\ur{تصویر} (tasvīr) — a picture}

\label{lesson:UR-C515-tasvir}



\begin{quote}
Before the new one: say the Urdu for a thimble, then the Urdu for a fringe.
\end{quote}

\subsection*{You'll want to know: \ur{تصویر}}
\textbf{\ur{تصویر}} — \emph{tasvīr} — \textquotedblleft{}a picture\textquotedblright{}.

\begin{quote}
\ur{دیوار} \ur{پر} \ur{دادا} \ur{کی} \ur{ایک} \ur{پرانی} \ur{تصویر} \ur{ہے۔}
\end{quote}

\begin{quote}
\emph{dīvār par dādā kī ek purānī tasvīr hai.}
\end{quote}

\begin{quote}
There is an old picture of grandfather on the wall.
\end{quote}

Feminine.

\begin{cousinweb}
\textbf{\ur{تصویر}} (\emph{tasvīr}) is usually traced, through Persian, to Arabic \emph{taṣwīr}, \textquotedblleft{}a picturing, a portrait\textquotedblright{}.
\end{cousinweb}

\begin{grammarlens}[title={What you've built}]
One more word for describing what you lived and felt.
\end{grammarlens}

\subsection*{Your turn}
\begin{itemize}
\raggedright
  \item \emph{Say it:} \emph{tasvīr}
  \item \emph{Say it:} \emph{tasvīr}, once more
  \item \emph{Say it:} \emph{dīvār par dādā kī ek purānī tasvīr hai.} — There is an old picture of grandfather on the wall.
  \item \emph{Recall:} say the Urdu for a thimble, then the Urdu for a fringe, then say \emph{tasvīr} again
\end{itemize}

\begin{tcolorbox}[breakable,colback=teal!4,colframe=teal!35!black,title={Before you move on}]
What does \ur{تصویر} mean? (\textquotedblleft{}A picture\textquotedblright{}.) Say it once more.
\end{tcolorbox}
\section[\texorpdfstring{\ur{مصور} (musavvir)}{مصور (musavvir)}]{\ur{مصور} (musavvir) — a painter (an artist)}

\label{lesson:UR-C515-musavvir}



\begin{quote}
Before the new one: say the Urdu for a fringe, then the Urdu for a picture.
\end{quote}

\subsection*{You'll want to know: \ur{مصور}}
\textbf{\ur{مصور}} — \emph{musavvir} — \textquotedblleft{}a painter (an artist)\textquotedblright{}.

\begin{quote}
\ur{اس} \ur{مصور} \ur{نے} \ur{پہاڑوں} \ur{کی} \ur{تصویریں} \ur{بنائیں۔}
\end{quote}

\begin{quote}
\emph{is musavvir ne pahāṛoṅ kī tasvīreṅ banāīṅ.}
\end{quote}

\begin{quote}
This painter made pictures of the mountains.
\end{quote}

Masculine or feminine, by the person meant.

\begin{cousinweb}
\textbf{\ur{مصور}} (\emph{musavvir}) is usually traced, through Persian, to Arabic \emph{muṣawwir}, \textquotedblleft{}one who gives form\textquotedblright{}.
\end{cousinweb}

\begin{grammarlens}[title={What you've built}]
One more word for describing what you lived and felt.
\end{grammarlens}

\subsection*{Your turn}
\begin{itemize}
\raggedright
  \item \emph{Say it:} \emph{musavvir}
  \item \emph{Say it:} \emph{musavvir}, once more
  \item \emph{Say it:} \emph{is musavvir ne pahāṛoṅ kī tasvīreṅ banāīṅ.} — This painter made pictures of the mountains.
  \item \emph{Recall:} say the Urdu for a fringe, then the Urdu for a picture, then say \emph{musavvir} again
\end{itemize}

\begin{tcolorbox}[breakable,colback=teal!4,colframe=teal!35!black,title={Before you move on}]
What does \ur{مصور} mean? (\textquotedblleft{}A painter (an artist)\textquotedblright{}.) Say it once more.
\end{tcolorbox}
\section[\texorpdfstring{\ur{خاکہ} (khākā)}{خاکہ (khākā)}]{\ur{خاکہ} (khākā) — a sketch; an outline}

\label{lesson:UR-C515-khaka}



\begin{quote}
Before the new one: say the Urdu for a picture, then the Urdu for a painter (an artist).
\end{quote}

\subsection*{You'll want to know: \ur{خاکہ}}
\textbf{\ur{خاکہ}} — \emph{khākā} — \textquotedblleft{}a sketch; an outline\textquotedblright{}.

\begin{quote}
\ur{مصور} \ur{نے} \ur{پہلے} \ur{پنسل} \ur{سے} \ur{خاکہ} \ur{بنایا۔}
\end{quote}

\begin{quote}
\emph{musavvir ne pehle pensil se khākā banāyā.}
\end{quote}

\begin{quote}
The painter first made a sketch in pencil.
\end{quote}

Masculine.

\begin{cousinweb}
\textbf{\ur{خاکہ}} (\emph{khākā}) is usually traced to Persian \emph{khāka}, \textquotedblleft{}a sketch, an outline\textquotedblright{}.
\end{cousinweb}

\begin{grammarlens}[title={What you've built}]
One more word for describing what you lived and felt.
\end{grammarlens}

\subsection*{Your turn}
\begin{itemize}
\raggedright
  \item \emph{Say it:} \emph{khākā}
  \item \emph{Say it:} \emph{khākā}, once more
  \item \emph{Say it:} \emph{musavvir ne pehle pensil se khākā banāyā.} — The painter first made a sketch in pencil.
  \item \emph{Recall:} say the Urdu for a picture, then the Urdu for a painter (an artist), then say \emph{khākā} again
\end{itemize}

\begin{tcolorbox}[breakable,colback=teal!4,colframe=teal!35!black,title={Before you move on}]
What does \ur{خاکہ} mean? (\textquotedblleft{}A sketch; an outline\textquotedblright{}.) Say it once more.
\end{tcolorbox}
\section[\texorpdfstring{\ur{شاہکار} (shāhkār)}{شاہکار (shāhkār)}]{\ur{شاہکار} (shāhkār) — a masterpiece}

\label{lesson:UR-C515-shahkar}



\begin{quote}
Before the new one: say the Urdu for a painter (an artist), then the Urdu for a sketch.
\end{quote}

\subsection*{You'll want to know: \ur{شاہکار}}
\textbf{\ur{شاہکار}} — \emph{shāhkār} — \textquotedblleft{}a masterpiece\textquotedblright{}.

\begin{quote}
\ur{یہ} \ur{تصویر} \ur{اس} \ur{مصور} \ur{کا} \ur{شاہکار} \ur{ہے۔}
\end{quote}

\begin{quote}
\emph{yih tasvīr is musavvir kā shāhkār hai.}
\end{quote}

\begin{quote}
This picture is the painter's masterpiece.
\end{quote}

Masculine.

\begin{cousinweb}
\textbf{\ur{شاہکار}} (\emph{shāhkār}) is usually traced to Persian \emph{shāhkār}, built from \emph{shāh}, \textquotedblleft{}a king\textquotedblright{}, and \emph{kār}, \textquotedblleft{}work\textquotedblright{}: a work fit for a king.
\end{cousinweb}

\begin{grammarlens}[title={What you've built}]
One more word for describing what you lived and felt.
\end{grammarlens}

\subsection*{Your turn}
\begin{itemize}
\raggedright
  \item \emph{Say it:} \emph{shāhkār}
  \item \emph{Say it:} \emph{shāhkār}, once more
  \item \emph{Say it:} \emph{yih tasvīr is musavvir kā shāhkār hai.} — This picture is the painter's masterpiece.
  \item \emph{Recall:} say the Urdu for a painter (an artist), then the Urdu for a sketch, then say \emph{shāhkār} again
\end{itemize}

\begin{tcolorbox}[breakable,colback=teal!4,colframe=teal!35!black,title={Before you move on}]
What does \ur{شاہکار} mean? (\textquotedblleft{}A masterpiece\textquotedblright{}.) Say it once more.
\end{tcolorbox}
\section[\texorpdfstring{\ur{جھلکنا} (jhalaknā)}{جھلکنا (jhalaknā)}]{\ur{جھلکنا} (jhalaknā) — to show through; to glint}

\label{lesson:UR-C515-jhalakna}



\begin{quote}
Before the new one: say the Urdu for a sketch, then the Urdu for a masterpiece.
\end{quote}

\subsection*{You'll want to know: \ur{جھلکنا}}
\textbf{\ur{جھلکنا}} — \emph{jhalaknā} — \textquotedblleft{}to show through; to glint\textquotedblright{}.

\begin{quote}
\ur{پانی} \ur{میں} \ur{چاند} \ur{جھلکتا} \ur{ہے۔}
\end{quote}

\begin{quote}
\emph{pānī meṅ chā̃d jhalaktā hai.}
\end{quote}

\begin{quote}
The moon glints in the water.
\end{quote}

Past (it …): \ur{وہ} \ur{چاند} \ur{پانی} \ur{میں} \ur{جھلکا}, \emph{voh chā̃d pānī meṅ jhalkā}.

\begin{cousinweb}
\textbf{\ur{جھلکنا}} (\emph{jhalaknā}) is inherited: CDIAL lists the Hindi cousin \emph{jhalaknā}, \textquotedblleft{}to sparkle\textquotedblright{}, under \emph{*jhal}, \textquotedblleft{}flash\textquotedblright{}, a form it reconstructs (CDIAL 5352).
\end{cousinweb}

\begin{grammarlens}[title={What you've built}]
One more word for describing what you lived and felt.
\end{grammarlens}

\subsection*{Your turn}
\begin{itemize}
\raggedright
  \item \emph{Say it:} \emph{jhalaknā}
  \item \emph{Say it:} \emph{jhalaknā}, once more
  \item \emph{Say it:} \emph{pānī meṅ chā̃d jhalaktā hai.} — The moon glints in the water.
  \item \emph{Recall:} say the Urdu for a sketch, then the Urdu for a masterpiece, then say \emph{jhalaknā} again
\end{itemize}

\begin{tcolorbox}[breakable,colback=teal!4,colframe=teal!35!black,title={Before you move on}]
What does \ur{جھلکنا} mean? (\textquotedblleft{}To show through; to glint\textquotedblright{}.) Say it once more.
\end{tcolorbox}
